\begin{abstract}

Multi-teacher on-policy distillation (MOPD) trains one student on its own
rollouts, scored token by token by a specialist teacher for each task. In
current methods, the task mixture can determine both the objective weights and
the amount of rollout and teacher computation assigned to each task. We fix the
objective weights and adapt only the number of fixed-size micro-batches assigned
to each task. Every teacher contributes to every optimizer step. Scaling each
micro-batch loss by its task weight divided by its count preserves the expected
gradient under every allocation, so the allocation changes gradient variance
without changing the objective. Gradient-Preconditioned Adaptive Sampling
(GPAS) assigns micro-batches in proportion to each task's objective weight and
gradient standard deviation after AdamW's diagonal scaling. Cost-GPAS also uses
the measured time per micro-batch and the fixed time per optimizer step. Both
methods use gradients and timings already recorded during training and require
no additional forward or backward passes. In a controlled calculation, GPAS
reduces scaled gradient variance to $0.679\times$ that of Uniform, whereas
allocation based on raw gradient noise increases it to $2.316\times$ that of
Uniform. We evaluate GPAS in four-teacher distillation of Qwen3-1.7B on
mathematics, code, instruction following, and science.
\missing{Insert the four-teacher loss, GPU-hour, and downstream-capability
results after the runs finish.}

\end{abstract}
